\section{Geometric inputs and their weighted forms}\label{sec:inputs}

The local arguments use three geometric estimates. The weighted form of
the Kakeya set theorem controls multiplicity for sufficiently dense
shadings. Multilinear Kakeya controls incidences with three transverse
directions. The planar Furstenberg theorem bounds the union of shadings
after projection onto a plane. We state these inputs and derive
the weighted forms needed for the indexed configurations of
Section~\ref{sec:critical}; the projection consequence is developed in
Section~\ref{sec:projections}.

Weights and shadings remain attached to individual indices, including
indices with the same trace. All constants associated with a bounded
coordinate chart are fixed independently of resolution. We write
$\eta_{\mathrm{in}}$ for the tolerances in the input theorems below;
these are unrelated to the fixed smoothing width $\eta$ in the temporal
cost.

For a bounded set $E\subset\R^j$, let $\mathcal N_\delta(E)$ be the
number of half-open cubes of the fixed mesh $\delta\mathbb Z^j$ that meet
$E$. Replacing this mesh by a translated mesh or using balls of
comparable radius changes covering numbers by a dimensional constant.
We say that a probability measure $\mu$ is
\emph{$(s,C)$-Frostman down to $\delta$} if
\begin{equation}\label{eq:input-frostman}
 \mu(B(x,r))\le Cr^s\qquad(\delta\le r\le1).
\end{equation}
For a finite separated set this refers to its uniform probability
measure. On bounded line charts we use the Euclidean distance between
slope--intercept parameters. The constants in
\eqref{eq:input-frostman} may change by a fixed factor under the bounded
coordinate changes used below.

\subsection{The set estimate and convex clustering}

A $\delta$-tube is the $\delta$-neighbourhood of a line segment of
length one. Changing between this convention and the cylinders in the
introduction costs only fixed changes of length and radius.
A shading is a measurable subset $Y(T)\subset T$.
For a finite set of such tubes put
\[
 \Delta_{\max}(\mathbb T)
 =\sup_{K\ \mathrm{convex}}
   \frac{\sum_{T\in\mathbb T:T\subset K}|T|}{|K|},
 \qquad
 \lambda(\mathbb T,Y)
 =\frac{\sum_{T\in\mathbb T}|Y(T)|}
        {\sum_{T\in\mathbb T}|T|}.
\]
Only convex bodies of positive volume containing a tube matter in
the supremum.

\begin{theorem}[Kakeya set estimate]\label{thm:set-input}
For every $\varepsilon>0$ there are $\eta_{\mathrm{in}}>0$ and $\delta_0>0$ such
that, if $0<\delta<\delta_0$,
\[
 \Delta_{\max}(\mathbb T)\le\delta^{-\eta_{\mathrm{in}}},
 \qquad \lambda(\mathbb T,Y)\ge\delta^{\eta_{\mathrm{in}}},
\]
then
\begin{equation}\label{eq:set-input}
 \left|\bigcup_{T\in\mathbb T}Y(T)\right|
 \ge\delta^\varepsilon\sum_{T\in\mathbb T}|T|.
\end{equation}
\end{theorem}

This is \citet[Theorem~1.1]{GWZ2026}. The hypothesis concerns convex
clustering, rather than separation of directions. The density
threshold is part of the statement and will be retained in every use.

We first extend this estimate to weighted indexed families. In particular,
two indices with the same tube may have different shadings.

\begin{lemma}[Weighted convex clustering]\label{lem:weighted-convex-set}
Fix $R\ge1$. For every $\varepsilon>0$ there are
$\eta_{\mathrm{in}}>0$ and $C<\infty$ with the following property. Let $(T_i)_i$
be finitely many $\delta$-tubes in $B(0,R)$, where $0<\delta<1$,
with repetitions allowed. Let $\omega_i>0$ and let
$Y_i\subset T_i$ be measurable shadings. Suppose
\begin{equation}\label{eq:weighted-convex-hyp}
 \sum_{i:T_i\subset K}\omega_i|T_i|\le A|K|
 \quad\hbox{for every convex body }K,
 \qquad |Y_i|\ge\delta^{\eta_{\mathrm{in}}}|T_i|.
\end{equation}
Then
\begin{equation}\label{eq:weighted-convex-conclusion}
 \sum_i\omega_i|Y_i|
 \le C A\delta^{-\varepsilon}\left|\bigcup_iY_i\right|.
\end{equation}
The constants are independent of the number of indices and of their
weights.
\end{lemma}

\begin{proof}
We sample the indices according to their normalized weights. A finite
family of box tests will control all convex clusters in the sample,
with only a logarithmic loss.

Write $v_\delta$ for the common tube volume, and normalize the
weights by $v_i=\omega_i/A$. Testing the convex body $T_i$ gives
$v_i\le1$. Testing a fixed ball containing the family gives
\[
 S:=\sum_i v_i\le C_R\delta^{-2}.
\]
Choose each index independently with probability $v_i$. The expected
number of chosen indices is $S$.

We first reduce all convex tests to finitely many box tests. If a
convex body $K$ contains a tube from the family, put
$K_R=K\cap\overline{B(0,R)}$; this intersection contains every tube
counted in $K$. Enlarge an inscribed John ellipsoid of $K_R$ by a
dimensional factor to obtain an ellipsoid containing $K_R$, with volume
at most a dimensional constant times $|K|$, largest axis $O_R(1)$,
and every axis at least a constant times $\delta$: the body contains
a ball of radius $\delta$. An enclosing rectangular box has the same
properties. Quantize its center, axis lengths, and orthonormal frame
with accuracy $c_R\delta^2$, and enlarge each side by $C_R\delta^2$.
There is thus a collection $\mathcal B_\delta$ of at most
$C_R\delta^{-C_R}$ convex boxes such that every relevant $K_R$ is
contained in some $Q\in\mathcal B_\delta$ satisfying
\[
 |Q|\le C_R|K|,\qquad |Q|/v_\delta\ge c_R.
\]
The quantization error is smaller than a fixed fraction of the
shortest side; this proves both containment and the volume comparison.

For $Q\in\mathcal B_\delta$, the number $Z_Q$ of chosen indices
whose tubes lie in $Q$ is a sum of independent Bernoulli variables,
and \eqref{eq:weighted-convex-hyp} gives
\[
 \E Z_Q\le |Q|/v_\delta.
\]
The elementary exponential-moment bound
\[
 \Prob(Z_Q\ge t)
 \le \inf_{a>0}\exp\bigl((e^a-1)\E Z_Q-at\bigr)
\]
shows that, for a sufficiently large fixed $L_R$, simultaneously
\begin{equation}\label{eq:sample-box-bound}
 Z_Q\le L_R(1+\log(1/\delta))\,|Q|/v_\delta
 \qquad(Q\in\mathcal B_\delta)
\end{equation}
with probability at least $3/4$. Indeed each failure probability
can be made smaller than $\delta^{C_R+2}$, and a union bound applies.
If $S\ge L_R(1+\log(1/\delta))$, the Bernoulli lower-tail bound
\[
 \Prob(N_{\rm sample}<S/2)\le e^{-S/8}
\]
gives at least $S/2$ chosen indices with probability at least $3/4$,
after increasing $L_R$. This lower-tail estimate follows by applying
the exponential moment method to $e^{-aN_{\rm sample}}$.
Fix a sample with both properties.

By the box comparison, its convex clustering is
$O_R(1+\log(1/\delta))$. The multiplicity of any repeated tube is
bounded by the same quantity, since a box of volume $O(v_\delta)$
contains that tube. Keep one chosen index for each distinct tube.
The resulting set has at least
$c_RS/(1+\log(1/\delta))$ members, and each retained shading still
has the required minimum density. It is not necessary to identify
different shadings on repeated tubes.

Apply Theorem~\ref{thm:set-input} with target loss
$\varepsilon/2$, and choose $\eta_{\mathrm{in}}$ no larger than its density
tolerance. For all sufficiently small $\delta$, the logarithmic
clustering bound satisfies the theorem's hypothesis. Consequently
\[
 \left|\bigcup_iY_i\right|
 \ge \frac{c_RS v_\delta\delta^{\varepsilon/2}}
             {1+\log(1/\delta)}.
\]
Since $\sum_i\omega_i|Y_i|\le ASv_\delta$, this proves
\eqref{eq:weighted-convex-conclusion}. If
$S<L_R(1+\log(1/\delta))$, use instead the pointwise bound
$\sum_i\omega_i\one_{Y_i}\le AS$ and absorb the logarithm.
The bounded range of remaining $\delta$ is covered by
$\sum_i\omega_i\le C_RA\delta^{-2}$, with a larger constant.
\end{proof}

\subsection{Full-time planks and indexed shadings}

The critical problem measures concentration in planks containing entire
affine traces. We now convert this plank condition into the convex
clustering condition of the preceding lemma, while counting only the
marked times in the shading.

Let $M_i(t)=b_i+tu_i\in\R^2$, $0\le t\le1$, where $b_i,u_i$
lie in a fixed bounded box. Put $\delta=D^{-1}$, where $D\ge1$.
A full-time plank
of widths $a,b$ in cell units is a set
\begin{equation}\label{eq:input-plank}
 \left\{(t,x):0\le t\le1,
   |\langle x-c-tv,e_1\rangle|\le a\delta,
   |\langle x-c-tv,e_2\rangle|\le b\delta\right\},
\end{equation}
where $(e_1,e_2)$ is an orthonormal spatial frame and $c,v\in\R^2$.
An index belongs to the plank when its entire trace belongs to it.

\begin{lemma}[Weighted set estimate in a line chart]\label{lem:weighted-set}
Fix a bounded line chart and a fixed bound for the permitted
enlargement of mesh cells. For every $\varepsilon>0$ there are
$\eta_{\mathrm{in}}>0$ and $C<\infty$ as follows. Suppose the indices have weights
$\omega_i>0$, and every full-time plank satisfies
\begin{equation}\label{eq:input-plank-bound}
 \sum_{i:M_i\subset P}\omega_i\le Aab
 \qquad(1\le a\le b\le D).
\end{equation}
Mark a set $J_i$ of
time bins of length $D^{-1}$ on each index, with
$\#J_i\ge D^{1-\eta_{\mathrm{in}}}$. Let $E$ be the union of the mesh cells
visited at those bin centers, or their fixed enlargements. Then
\begin{equation}\label{eq:weighted-chart-conclusion}
 \sum_i\omega_i\#J_i\le C A D^\varepsilon\#E.
\end{equation}
In particular, when every index has $k$ marked bins and
$n=\sum_i\omega_i$, its raw average multiplicity satisfies
$nk/\#E\le C A D^\varepsilon$.
\end{lemma}

The plank bound also holds, up to a fixed multiplicative change in
$A$, for fixed dilations and widths between fixed multiples of $1$
and $D$. This follows from the endpoint partition in the proof and
is not an additional hypothesis.

\begin{proof}
There are two geometric steps: a plank bound controls weighted convex
clustering of thickened traces, and disjoint balls at marked bin centers
convert incidence counts to shading volume.

It is harmless to allow $a,b$ up to a fixed multiple of $D$ in
\eqref{eq:input-plank-bound}. A larger spatial coordinate can be
handled by partitioning its two endpoint values, at $t=0$ and
$t=1$, into intervals of fixed length in physical units.
Linear interpolation puts every part inside a permitted
full-time plank. There are only a bounded number of parts,
depending on the chart. Apply the given bound to each part;
the sum of their area factors is at most a fixed multiple
of $ab$. A fixed decrease in the lower width cutoff is handled
by enlargement, at a fixed cost in the area factor.

Enclose each trace in a tube $T_i$ of radius $C_0\delta$ and a
common fixed Euclidean length $L$. Choose $L$ large enough that the
trace, and balls of radius $C_0\delta$ about its points, lie away from the
end caps. Both constants depend only on the bounded chart.

We check weighted convex clustering for these tubes. Let $K$
contain some of them, intersect $K$ with a fixed ambient ball
containing all $T_i$, and take an enclosing ellipsoid $\mathcal E$
of volume $O(|K|)$ by dilating an inscribed John ellipsoid.
Its diameter is bounded. Its horizontal
slices have affine centers and a fixed pair of principal axes; their
semiaxes are a common scalar multiple of those of the maximal slice.
The time diameter of $\mathcal E$ is at least one, because every
counted tube contains its full chart trace. The ellipsoid volume is
a dimensional constant times this time diameter times the maximal
slice area. Thus that area is $O(|K|)$.

Circumscribe a rectangle about the maximal horizontal slice and
translate it to the affine center of each slice. The resulting full-time
plank contains all the counted traces and has cross-sectional area at
most $C|K|$.
Its spatial widths are
at least a constant times $\delta$, since a counted tube contains
a spatial disk of comparable radius at every chart time.
Applying \eqref{eq:input-plank-bound} therefore gives
\[
 \sum_{i:T_i\subset K}\omega_i\le C A\delta^{-2}|K|.
\]
Since $|T_i|\asymp\delta^2$, the weighted convex clustering of the
family is at most $CA$.

At each marked bin center put a ball of radius $c_0\delta$ on the
trace, with $c_0>0$ fixed and small. Balls belonging to different
bins of one index are disjoint. Their union $Y_i\subset T_i$ has
volume comparable to $\#J_i\delta^3$; its density in $T_i$ is at
least $c\delta^{\eta_{\mathrm{in}}}$. Moreover
\[
 \left|\bigcup_iY_i\right|\le C\delta^3\#E.
\]
Rescale the common length $L$ to one and apply
Lemma~\ref{lem:weighted-convex-set}. Choose the present $\eta_{\mathrm{in}}$
smaller than half the density tolerance of that lemma, so fixed
length and density constants are absorbed for small $\delta$.
Canceling $\delta^3$ proves \eqref{eq:weighted-chart-conclusion}
for large $D$. For bounded $D$, a fixed collection of widest
planks bounds the total weight by $CA D^2$. The pointwise
multiplicity bound by this total weight gives the conclusion
after increasing $C$ on that bounded range.
\end{proof}

\subsection{The planar Furstenberg estimate}

The planar input will be applied to the neighbors of a reference line
in a nearly coplanar configuration. Each neighbor keeps its marked times,
so we need a lower bound that retains the number of shading cells as a
factor. The following dual form of the quantitative planar theorem does
so.

\begin{theorem}[Quantitative planar Furstenberg estimate]
\label{thm:furstenberg-input}
Fix $0<s\le1$, $0<t\le2$, and a bounded line chart. For every
$\varepsilon>0$ there are $\eta_{\mathrm{in}}>0$, $c>0$, and $\delta_0>0$ with
the following property. Let $\mathcal L$ be a finite
$\delta$-separated family of lines whose uniform parameter
distribution is $(t,\delta^{-\eta_{\mathrm{in}}})$-Frostman down to $\delta$.
For every $\lambda\in\mathcal L$, let $Y(\lambda)$ be a
$\delta$-separated subset of a fixed bounded part of $\lambda$,
with
\[
 k\le\#Y(\lambda)\le2k,
\]
whose uniform distribution is
$(s,\delta^{-\eta_{\mathrm{in}}})$-Frostman down to $\delta$. Then, for
$0<\delta<\delta_0$,
\begin{equation}\label{eq:furstenberg-input}
 \mathcal N_\delta\left(\bigcup_{\lambda\in\mathcal L}Y(\lambda)\right)
 \ge c k\delta^{-\min\{t,(s+t)/2,1\}+\varepsilon}.
\end{equation}
The same assertion holds for mesh-cell shadings within a fixed
multiple of $\delta$ of the lines, with fixed changes of constants.
\end{theorem}

This follows from \citet[Theorem~4.1 and Definitions~1.3, 2.3, 3.1]{RW2025}
by point--line duality. To spell out the translation, in coordinates
where $\lambda$ has equation $y=ax+b$, it corresponds to the
parameter point $(a,b)$, and a shading point $(x,y)$ corresponds to
the dual line $b=y-ax$. Incidence is preserved. Along a bounded-slope
original line, the shading coordinate $x$ is bi-Lipschitz to
arclength, so its Frostman bound is exactly the required bound for
the incident dual directions. Discretizing introduces only fixed
changes of cell size and bounded overlaps. In fact, a parameter
cell contains only a bounded number of the separated original
lines, and a dual parameter cell contains only a bounded number
of separated shading points from any one original line. Retaining
representatives preserves relative Frostman counts up to fixed
factors. The nice configuration of Ren--Wang therefore has
$M\asymp k$ selected incident tubes at each base cell; the number
of distinct dual parameter cells is comparable to the covering
number in \eqref{eq:furstenberg-input}. Choose the displayed
$\eta_{\mathrm{in}}$ smaller than half the source theorem's tolerance to absorb
the fixed factors. A finite division into coordinate
charts handles all bounded planar patches. Their theorem permits
the constant $\delta^{-\eta_{\mathrm{in}}}$; this quantitative tolerance is why
subpower losses are allowed here.

The cited theorem is stated at dyadic scales. A comparable dyadic
scale changes mesh sizes and covering numbers only by fixed factors,
which are absorbed in the displayed constants.

The equivalent shaded-tube formulation in
\citet[Theorem~0.6, Equation~(0.8)]{WW2024} gives the same exponent.
Indeed its density parameter is $\lambda\asymp k\delta$; division
of the area bound by $\delta^2$ gives the factor $k$ in
\eqref{eq:furstenberg-input}.

\begin{corollary}[Weighted planar families]\label{cor:weighted-furstenberg}
Fix $0<s\le1$, $0<t\le2$, and a bounded line chart. For every
$\varepsilon>0$ there are $\eta_{\mathrm{in}}>0$, $c>0$, and $\delta_0>0$ as follows.
Let $0<\delta<\delta_0$ and let finitely many, possibly repeated,
line indices $i$ carry positive weights $\omega_i$. Suppose their
normalized line-parameter measure is
$(t,\delta^{-\eta_{\mathrm{in}}})$-Frostman down to $\delta$. For each index $i$,
let $Y_i$ be a $\delta$-separated shading of its line in a fixed
bounded chart, with $k\le\#Y_i\le2k$ and with its uniform
distribution $(s,\delta^{-\eta_{\mathrm{in}}})$-Frostman down to $\delta$. Then
\[
 \mathcal N_\delta\left(\bigcup_iY_i\right)
 \ge c k\delta^{-\min\{t,(s+t)/2,1\}+\varepsilon}.
\]
The constants are independent of the number of indices, their
repetitions, and their weights.
\end{corollary}

\begin{proof}
It suffices to select separated line parameters whose uniform
distribution inherits the weighted Frostman bound, and to keep one
original shading at each selected parameter.

Push the normalized weights to the $O(\delta^{-2})$ parameter cells.
Discard cells of mass below $\delta^4$; for small $\delta$ the lost
mass is less than $1/2$. There are $O(1+\log(1/\delta))$ remaining
dyadic mass levels, so one level has mass at least
$c/(1+\log(1/\delta))$. On this level the uniform distribution of
occupied cells has Frostman constant at most
$C\delta^{-\eta_{\mathrm{in}}}(1+\log(1/\delta))$: compare the comparable masses
of the cells in an enlarged ball with the original probability.
Color the parameter mesh by $J$ congruence classes, where $J$ is
a fixed constant, and retain a class containing at least $1/J$
of the retained cells. Choose one actual line
and its original shading from each retained cell. The resulting
line parameters are separated, the shadings retain their Frostman
bounds, and their union is contained in the original union.
Apply Theorem~\ref{thm:furstenberg-input}, having chosen the current
$\eta_{\mathrm{in}}$ smaller than half its allowed tolerance.
\end{proof}

\begin{remark}\label{rem:input-subpower-shadings}
The same conclusion in normalized-log estimates holds when the
shading cardinalities are comparable only up to subpowers of
$\delta^{-1}$. Thin every shading to the common minimum cardinality.
Its relative Frostman constant increases by at most the ratio of
the old and new cardinalities. That ratio, and the changes just
used to regularize line weights, are subpowers and are absorbed by
the fixed positive tolerance in Theorem~\ref{thm:furstenberg-input}.
The case $t<s$ is included: the minimum in
\eqref{eq:furstenberg-input} is then $t$.
\end{remark}

\subsection{Multilinear Kakeya with quantitative transversality}

The remaining input limits overlap among three families with transverse
directions. We retain its dependence on the determinant of those
directions, since the local arguments allow that determinant to decrease
with the resolution.

\begin{theorem}[Endpoint multilinear Kakeya]\label{thm:multilinear-input}
For $j=1,2,3$, let $T_{j,a}$ be finitely many infinite cylinders of
radius one in $\R^3$, with unit axial directions $v_{j,a}$. Suppose
\[
 |\det(v_{1,a_1},v_{2,a_2},v_{3,a_3})|\ge\theta>0
 \quad\hbox{for every }(a_1,a_2,a_3).
\]
For nonnegative weights $\omega_{j,a}$ put
$\mu_j=\sum_a\omega_{j,a}\one_{T_{j,a}}$ and
$M_j=\sum_a\omega_{j,a}$. Then
\begin{equation}\label{eq:weighted-multilinear}
 \int_{\R^3}(\mu_1\mu_2\mu_3)^{1/2}
 \le C\theta^{-1/2}(M_1M_2M_3)^{1/2}.
\end{equation}
In particular this holds for finite length-$D$ tubes with no
additional dependence on $D$.
\end{theorem}

\begin{proof}
For unit weights this is
\citet[Theorem~1.3, p.~264]{Guth2010}, the endpoint form of
multilinear Kakeya introduced in
\citet[Theorem~1.15]{BCT2006}. Integer weights are represented by
repetition. There is no distinctness requirement; alternatively
one can first separate repeated cylinders by distinct parallel
translations tending to zero and use Fatou's lemma. Off their
boundaries, the indicators then converge to the repeated indicators,
and directions are unchanged.

For rational weights clear denominators separately in the three
families; the factors cancel on the two sides of
\eqref{eq:weighted-multilinear}. Approximate real weights
monotonically from below by rational weights. Monotone convergence
proves the displayed estimate. Truncating the cylinders can only
decrease the integrand.
\end{proof}

\begin{corollary}[Cube counts and small caps]\label{cor:multilinear-cubes}
For unit mesh cubes $Q$, let
$m_j(Q)=\sum_{a:T_{j,a}\cap Q\ne\varnothing}\omega_{j,a}$.
Under the hypotheses of Theorem~\ref{thm:multilinear-input},
\begin{equation}\label{eq:multilinear-cubes}
 \sum_Q(m_1(Q)m_2(Q)m_3(Q))^{1/2}
 \le C\theta^{-1/2}(M_1M_2M_3)^{1/2}.
\end{equation}
If the directions lie in caps of radius at most $c\theta$ about
centers whose determinant is at least $\theta$, the same conclusion
holds with another absolute constant, for sufficiently small $c$.
\end{corollary}

\begin{proof}
Increase each cylinder's radius to $1+\sqrt3$. Every cube meeting
the original cylinder lies entirely in the enlarged cylinder, so
the enlarged multiplicity dominates $m_j(Q)$ everywhere on $Q$.
Integrate over the disjoint cubes and use
Theorem~\ref{thm:multilinear-input}, followed by the fixed radius
rescaling. Finally the determinant of three unit vectors changes
by at most a constant times the sum of their perturbations.
Caps of radius $c\theta$ therefore give cross-family determinant
at least $\theta/2$ when $c$ is small enough.
\end{proof}

For later use, a cap partition of mesh $r$ has $O(r^{-2})$ members.
Summing \eqref{eq:multilinear-cubes} over triples of caps costs at
most $O(r^{-6})$. Thus choosing $r$ to be a fixed positive power
of $\theta$ introduces only a polynomial loss in $\theta^{-1}$;
in particular the choice $r=\theta^2$ is permissible as
$\theta\to0$.
